%%%%%%%%%%%%%%%%%%%%%%%%%%%%%%%%%%%%%%%%%%%%%%%%%%%%%%%%%%%%%%
%%%%%%%%%%%% CA-0411 Análisis de Datos I (propuesta) %%%%%%%%%%
%%%%%%%%%%%%%%%%%%%%%%%%%%%%%%%%%%%%%%%%%%%%%%%%%%%%%%%%%%%%%%
% Propuesta a evaluar por separado. Fuentes resumidas:
% - Cartas al estudiantado CA-0411 (I-2024, I-2025, I-2026),
%   https://emate.ucr.ac.cr/carrera-ciencias-actuariales
% - Programa vigente CA-0411 Análisis de Datos I (Cursos/)
% No modificar sin decisión curricular expresa.
\newpage
\subsection*{CA-0411 Análisis de Datos I}
\addcontentsline{toc}{subsection}{CA-0411 Análisis de Datos I (propuesta)}

\hypertarget{CA0411AnalisisDatosPropuesta}{}
\begin{refsection}
\begin{table}[H]
\centering{
\begin{tabular}{|ll|}
\hline
\textbf{Año:} IV  & \textbf{Requisitos:} CA-0303 Estadística I \\
\textbf{Ciclo:} VII  & \textbf{Correquisitos:} Ninguno \\
\textbf{Tipo de curso:} Teórico-práctico & \textbf{Horas semanales:} 5 \\
\textbf{Créditos:} 4 & \textbf{Virtualidad:} P, PV \\ \hline
\end{tabular}
}
\end{table}

\subsubsection*{Descripción del curso}

Este curso está dirigido a estudiantes avanzados que han completado una formación básica en estadística e inferencia, por lo que se ubica en el cuarto año (ciclo~VII) de la ruta de matemática aplicada. Introduce la modelación de aprendizaje estadístico y la implementación de métodos de aprendizaje no supervisado y supervisado, de modo que la persona estudiante adquiera herramientas de análisis de datos con aplicaciones en múltiples disciplinas, incluyendo riesgos, seguros, finanzas, demografía y otras ciencias cuantitativas.

La propuesta resume y adapta, al formato de programas de la carrera de Matemática, los contenidos y énfasis de las cartas al estudiantado de \textbf{CA-0411 Análisis de Datos~I} (ofertas recientes) y el programa existente de análisis de datos~I.

\subsubsection*{Objetivo general}

Fortalecer en el estudiantado las competencias teóricas y analíticas del aprendizaje estadístico para diseñar, implementar y evaluar críticamente modelos de aprendizaje supervisado y no supervisado aplicados a datos reales.

\subsubsection*{Objetivos específicos}

Al finalizar el curso el/la estudiante estará en la capacidad de:

\begin{enumerate}
\item Comprender el modelo de aprendizaje estadístico y su relación con el proceso de modelación en ciencia de datos. (SC11, SH05.)
\item Implementar modelos de aprendizaje no supervisado y supervisado en conjuntos de datos de distinta índole. (SC11, SH05.)
\item Analizar los procesos de construcción, validación y calibración de modelos, con especial énfasis en aplicaciones cuantitativas (por ejemplo, riesgos, seguros, finanzas o demografía). (SC11, SH05.)
\item Comunicar resultados de modelación a públicos especializados y no especializados, mediante informes y presentaciones. (SH09.)
\end{enumerate}

\subsubsection*{Contenidos}

La distribución por semanas es orientativa. La persona docente puede reordenar la práctica y profundizar subtemas, sin omitir los contenidos obligatorios. Los ítems marcados como \emph{ampliación} pueden cubrirse según el perfil del grupo y el tiempo disponible.

\begin{enumerate}
\item \textbf{Introducción al aprendizaje estadístico y al análisis multidimensional (2 semanas):}
\begin{enumerate}
\item Motivación del aprendizaje estadístico; marco de modelación en ciencia de datos; sesgo--varianza; conjuntos de entrenamiento, validación y prueba.
\item Tabla de datos; nubes de individuos y nubes de variables.
\item Métrica en un espacio euclidiano y distancia en un espacio afín; métrica de pesos.
\item Centro de gravedad de la nube; matriz de varianza--covarianza; inercias.
\end{enumerate}
\item \textbf{Aprendizaje no supervisado (5--6 semanas):}
\begin{enumerate}
\item \textbf{Análisis en componentes principales (ACP):}
\begin{enumerate}
\item Aplicación e interpretación de un ACP.
\item Algoritmo e implementación del ACP.
\item Presentación geométrica y algebraica del ACP.
\end{enumerate}
\item \textbf{Clasificación automática (conglomerados):}
\begin{enumerate}
\item Panorama de algoritmos de clasificación automática.
\item Método de $k$-medias.
\item Método de $k$-medoïdes.
\item Agrupación (clasificación) jerárquica.
\end{enumerate}
\item \textbf{Ampliaciones de análisis multidimensional} (a criterio docente):
\begin{enumerate}
\item Escalamiento multidimensional (\emph{multidimensional scaling}).
\item Análisis de correspondencias simples (AFC).
\item Análisis de correspondencias múltiples (AFCM).
\end{enumerate}
\end{enumerate}
\item \textbf{Aprendizaje supervisado (6--7 semanas):}
\begin{enumerate}
\item Motivación del aprendizaje supervisado.
\item Método de $k$ vecinos más cercanos (KNN) (\emph{ampliación} o puente hacia evaluación de clasificadores).
\item \textbf{Medición de la calidad de un método supervisado:}
\begin{enumerate}
\item Matriz de confusión.
\item Análisis del error y curvas ROC.
\end{enumerate}
\item \textbf{Análisis discriminante:}
\begin{enumerate}
\item Enfoque factorial.
\item Enfoque bayesiano / métodos bayesianos asociados.
\end{enumerate}
\item Árboles de decisión (método CART).
\item Bosques aleatorios.
\item Valores de Shapley para explicar la predicción de los modelos.
\item Máquinas de soporte vectorial.
\item Validación cruzada y calibración de modelos.
\end{enumerate}
\item \textbf{Proyecto integrador} (transversal; énfasis en las semanas finales):
\begin{enumerate}
\item Planteamiento de una pregunta de investigación y uso de datos reales.
\item Aplicación de al menos un método de aprendizaje supervisado o no supervisado \emph{no desarrollado en detalle} en clase.
\item Bitácoras de avance, informe escrito y presentación oral; se recomienda orientar el tema, cuando los datos lo permitan, a riesgos, seguros, finanzas, crédito o segmentación.
\end{enumerate}
\end{enumerate}

\subsubsection*{Metodología}

Se espera que la persona docente aplique al menos tres de las siguientes estrategias didácticas:

\begin{enumerate}
\item Clases magistrales con casos prácticos de modelación.
\item Laboratorios guiados en R o Python (el curso puede impartirse preferentemente en uno de los dos lenguajes; se permite que el estudiantado trabaje en el otro si demuestra dominio equivalente).
\item Proyecto grupal a lo largo del ciclo, con bitácoras de avance.
\item Búsqueda de datos y preguntas de investigación en fuentes externas.
\item Solución de problemas reales ligados a la industria o a contextos interdisciplinarios.
\item Visitas o charlas de personas expertas (banca, riesgos, segmentación u otros ámbitos de valor agregado).
\item Identificación de conceptos en artículos científicos sencillos e interpretación de resultados.
\item Peer project learning y retroalimentación entre pares en presentaciones.
\end{enumerate}

Se motiva a la persona docente a incentivar un proyecto con alto componente computacional durante las 15~semanas del curso, la aplicación de modelos en la cuantificación de incertidumbre o riesgo, y la interpretación de resultados hacia público no especializado. Cuando la disponibilidad de datos lo permita, se recomienda orientar el proyecto a temas de riesgos, seguros, finanzas, crédito o segmentación, sin excluir otras aplicaciones matemáticas o científicas coherentes con el perfil del estudiantado.

\subsubsection*{Evaluación}

La evaluación del aprovechamiento puede combinar, a criterio de la persona docente, los siguientes rubros (con ponderaciones definidas en la carta al estudiantado de cada oferta):

\begin{enumerate}
\item Tareas teóricas y de implementación.
\item Laboratorios evaluados (componente teórico y de programación).
\item Proyecto: bitácoras, presentaciones orales e informe escrito.
\item Examen o prueba integradora (opcional según la oferta).
\item Documentos sintéticos: póster, resumen (\emph{abstract}), etc.
\end{enumerate}

\subsubsection*{Bibliografía}

\begin{enumerate}
\item James, G., Witten, D., Hastie, T., Tibshirani, R. y Taylor, J. (2023). \emph{An Introduction to Statistical Learning: With Applications in Python.} Springer. También la edición con aplicaciones en R.
\item Hastie, T., Tibshirani, R. y Friedman, J. (2009). \emph{The Elements of Statistical Learning: Data Mining, Inference, and Prediction} (2nd ed.). Springer.
\item Trejos, J.; Castillo, W.; González, J. (2014). \emph{Análisis Multivariado de Datos. Métodos y Aplicaciones.} Editorial de la Universidad de Costa Rica.
\item Breiman, L., Friedman, J., Olshen, R. y Stone, C. (1984). \emph{Classification and Regression Trees.} Chapman and Hall/CRC.
\item Breiman, L. (2001). Random forests. \emph{Machine Learning}, 45(1), 5--32.
\item Lundberg, S. M. y Lee, S.-I. (2017). A unified approach to interpreting model predictions. En \emph{Advances in Neural Information Processing Systems}.
\item Géron, A. (2022). \emph{Hands-On Machine Learning with Scikit-Learn, Keras, and TensorFlow.} O'Reilly.
\item Wickham, H. y Grolemund, G. (2017). \emph{R for Data Science.} O'Reilly.
\item McKinney, W. (2022). \emph{Python for Data Analysis.} O'Reilly.
\item Hair, J. F., Black, W. C., Babin, B. J. y Anderson, R. E. (2019). \emph{Multivariate Data Analysis} (8th ed.). Cengage.
\item Murphy, K. P. (2012). \emph{Machine Learning: A Probabilistic Perspective.} MIT Press.
\item Yu, B. y Barter, R. \emph{Veridical Data Science: The Practice of Responsible Data Analysis and Decision Making.} \url{https://vdsbook.com/}
\end{enumerate}

\end{refsection}
